	
\begin{chineseabstract}
本文档阐述了应用于 14-bit 高精度数据采集系统的模拟前端（AFE）及全差分电容数字模拟转换器
（CDAC）的电路架构选型与设计过程。面对低底噪、高线性度与大容性负载驱动等多方面的设计要求，
本设计构建了一套包含 S2D 缓冲级、带密勒补偿的高增益（$138\,\mathrm{dB}$）FDA 核心，
以及双 8-bit R-2R 梯形电阻网络的全差分信号调理链路。

在系统级噪声统筹中，结合 Friis 噪声级联公式，对前级施以较高偏置电流以压制闪烁噪声与热噪声平台
；同时，依托连续时间共模反馈（CMFB）环路，较好地抑制了电源波动与共模漂移对动态量程的影响。
在 CDAC 设计层面，引入了电容分段桥接技术与 $V_{\mathrm{cm}}$ 时序切换策略，
以期降低开关能耗与面积开销；此外，利用基于离散时间域线性回归的 CMOS 传输门宏模型，
对开关阵列的充放电建立过程进行了初步的数学建模与尺寸加权配置。本报告基本梳理了从系统指标分
解到晶体管级参数推演的设计流程。然而，宏模型在极端工艺偏差下的拟合精度、以及双 R-2R 网络在
实际流片中的绝对匹配度，仍有待后续硅片测试结果的进一步验证与校准。
\chinesekeyword{模拟前端设计；全差分架构；可编程增益放大器；R-2R网络；传输门宏模型；建立时间}

\end{chineseabstract}

